%Introducción al tema. Página 1


Buena parte de los sistemas de ingeniería electrónica o de
telecomunicaciones procesan información en forma de
\emph{señales}, que son sometidas a diversas operaciones, por
ejemplo:

\begin{itemize}
\item Transducción de una magnitud física a otra, por ej: $p(t)\rightarrow v(t)$ (sonido $\rightarrow$
voltaje).
\item Transmisión / Recepción (a través de un medio que transporta
ondas).
\item Conversión Análogo/Digital, Digital/Analógica.
\item Modulación / detección (para acondicionar una señal a un
canal).
\item Codificación  (para corregir errores o seguridad).
\end{itemize}
Idealmente, las operaciones en los ejemplos anteriores serían
todas \emph{reversibles}. En la práctica, la inversión es
imperfecta, por diferentes razones: distorsión, errores de
cuantificación, ruido.

También interesan operaciones no reversibles de ``procesamiento de
señales", por ejemplo filtrado,
compresión, segmentación.

En este curso veremos en detalle las operaciones de Modulación /
Detección digital.